\documentclass[10pt,a4paper]{article}

\usepackage[utf8]{inputenc}
\usepackage[T1]{fontenc}
\usepackage{amsmath,mathtools,amsfonts,amssymb}
\usepackage{graphicx}
\usepackage[spanish, es-tabla]{babel}
\usepackage{lastpage,tabto,xcolor}
\usepackage[a4paper, top=0.9in, bottom=1in, left=0.9in, right=0.9in]{geometry}
\usepackage{fancyhdr}
\usepackage{titling}
\usepackage{float}
\usepackage{enumitem}
\usepackage{booktabs}
\usepackage[mode=image]{standalone}

\usepackage{tikz}
\usepackage{pgfplots}
\pgfplotsset{compat=1.18}
\usepackage{caption}
\usepackage{framed}
\usetikzlibrary{arrows.meta, babel}

\definecolor{utnblue}{RGB}{0, 51, 102}
\definecolor{utnred}{RGB}{204, 0, 0}
\definecolor{utngreen}{RGB}{0, 153, 76}

% Fecha de versión automática según fecha de compilación
\newcommand{\verdate}{\the\year.\ifnum\month<10 0\fi\the\month.\ifnum\day<10 0\fi\the\day}

% Incisos con letras: a), b), c), ...
\setlist[enumerate,1]{label=\textbf{\alph*)}, leftmargin=*, itemsep=4pt}

\setlength{\headsep}{12pt}
\setlength{\droptitle}{-0.5in}
\pretitle{\begin{center}\vspace{-1em}\normalfont\Large}
\posttitle{\par\end{center}\vspace{-2.5em}}
\preauthor{\begin{center}\vspace{-0.5em}}
\postauthor{\par\end{center}}
\predate{}\postdate{}

\pagestyle{fancy}
\lhead{\footnotesize Análisis de Señales y Sistemas  - Año \the\year{} Ver. \verdate}
\rhead{\footnotesize Resolución del Trabajo Práctico N$^\text{o}$7}
\rfoot{\footnotesize Página \thepage\ de \pageref{LastPage}}
\renewcommand{\headrulewidth}{0.4pt}
\renewcommand{\footrulewidth}{0.4pt}

% Encabezado de tema
\newcommand{\tema}[1]{%
  \bigskip
  \section*{#1}
  \vspace{-0.6em}
}

% Encabezado de ejercicio
\newcommand{\ejercicio}[1]{\bigskip\noindent{\large\textbf{Ejercicio #1}}\par\nobreak\smallskip}

% Consigna transcripta del enunciado
\newenvironment{consigna}{%
  \par\smallskip
  \def\FrameCommand{\textcolor{utnblue!35}{\vrule width 2pt}\hspace{10pt}}%
  \MakeFramed{\advance\hsize-\width \FrameRestore}%
  \small\noindent\textit{Enunciado.}\enspace\ignorespaces
}{\endMakeFramed\par\smallskip}

\title{%
	{\normalsize Departamento de Ingeniería en Tecnologías Electrónicas, UTN FRM}\\[0.1cm]
	{\large Análisis de Señales y Sistemas}\\[0.15cm]
	{\normalsize Resolución del Trabajo Práctico N$^{\text o}$7: Sistemas de tiempo discreto}%
}
\author{}
\date{}

\begin{document}
	\maketitle
	\thispagestyle{fancy}

	\label{sec:resolucion_tp7}

	\noindent
	Antes de cada bloque de ejercicios va un repaso breve de las herramientas que se usan
	en él. La notación es la del trabajo anterior: las señales de tiempo continuo llevan
	paréntesis, $x(t)$, y las secuencias corchetes, $x[n]$, con $n$ entero. El operador
	$\mathcal{T}\{\cdot\}$ nombra la regla del sistema, $h[n]$ su respuesta al impulso y
	$*$ la suma de convolución. La notación por extensión
	$x[n]=\{\,\underset{\uparrow}{a} \quad b \quad c\,\}$ lista los valores de la
	secuencia con la flecha marcando el índice $n=0$, y fuera de esos índices la
	secuencia vale cero.

	% ============================================================================
	\tema{Propiedades de los sistemas discretos}
	% ============================================================================

	Las seis propiedades con que se clasifica un sistema discreto se aplican sobre su
	regla, sin necesidad de conocer ninguna entrada en particular.

	\begin{itemize}[nosep]
		\item \textbf{Memoria.} El sistema es \emph{sin memoria} si $y[n]$ depende
		      únicamente de $x[n]$. Cualquier dependencia de otras muestras de la entrada,
		      o de valores anteriores de la salida, lo vuelve un sistema \emph{con memoria}.
		\item \textbf{Causalidad.} Es \emph{causal} si $y[n]$ se construye solo con
		      muestras $x[k]$ de índice $k \le n$.
		\item \textbf{Estabilidad BIBO.} Es \emph{estable} si toda entrada acotada,
		      $|x[n]| \le M_x < \infty$, produce una salida acotada.
		\item \textbf{Invertibilidad.} Es \emph{invertible} si entradas distintas producen
		      salidas distintas.
		\item \textbf{Linealidad.}
		      $\mathcal{T}\{\alpha x_1[n] + \beta x_2[n]\}
		      = \alpha\,\mathcal{T}\{x_1[n]\} + \beta\,\mathcal{T}\{x_2[n]\}$.
		\item \textbf{Invarianza temporal.} Si $x[n]$ produce $y[n]$, entonces
		      $x[n-n_0]$ produce $y[n-n_0]$ para todo $n_0$ entero.
	\end{itemize}

	Las dos últimas se verifican siempre con la definición, y para negarlas alcanza con un
	único contraejemplo. Hay además una regla práctica que ahorra trabajo: si la regla
	combina desplazamientos de la entrada y de la salida con coeficientes
	\emph{constantes}, el sistema es lineal e invariante en el tiempo. La implicación
	corre en un solo sentido, de modo que cuando la regla no tiene esa forma no se
	concluye nada y hay que volver a la definición.

	% ----------------------------------------------------------------------------
	\ejercicio{1: clasificación de cuatro sistemas}

	\begin{consigna}
	Sean los sistemas
	\begin{align*}
		\mathcal{T}_1&: \; y[n] = x[n] - 2\,x[n-2], \\
		\mathcal{T}_2&: \; y[n] = x[n]\,x[n-1], \\
		\mathcal{T}_3&: \; y[n] = \textstyle\sum_{k=n-2}^{n+2} x[k], \\
		\mathcal{T}_4&: \; y[n] = (n+1)\,x[n].
	\end{align*}
	\begin{enumerate}
		\item Indicar cuáles tienen memoria y cuáles son causales.
		\item Determinar cuáles son lineales y cuáles invariantes en el
		tiempo, exhibiendo un contraejemplo en cada caso negativo.
		\item Determinar cuáles son estables BIBO.
		\item Señalar cuál de los cuatro no puede usarse en un
		procesamiento en tiempo real e indicar en qué situación sí es
		utilizable.
	\end{enumerate}
	\end{consigna}

	Las dos primeras propiedades salen de leer qué muestras usa cada regla, sin necesidad de
	ninguna cuenta. El sistema $\mathcal{T}_3$ es el promediador
	centrado sin el factor $1/5$, y su suma recorre los cinco índices de $n-2$ a $n+2$, es
	decir
	\[
		y[n] = x[n-2]+x[n-1]+x[n]+x[n+1]+x[n+2].
	\]

	\begin{enumerate}
		\item Solo $\mathcal{T}_4$ es sin memoria, porque su
		salida en el instante $n$ se arma con la muestra $x[n]$ y nada más. El factor
		$(n+1)$ no aporta memoria, ya que es un número que el sistema conoce por el índice
		en que está trabajando y no un valor guardado de la entrada. Los otros tres tienen
		memoria, $\mathcal{T}_1$ porque usa $x[n-2]$, $\mathcal{T}_2$ porque usa $x[n-1]$ y
		$\mathcal{T}_3$ porque usa cuatro muestras además de la actual.

		Para la causalidad hay que mirar si alguna regla usa índices mayores que $n$.
		En $\mathcal{T}_1$, $\mathcal{T}_2$ y $\mathcal{T}_4$ todos los índices son menores
		o iguales que $n$, de modo que los tres son causales. En $\mathcal{T}_3$ aparecen
		$x[n+1]$ y $x[n+2]$, dos muestras posteriores al instante que se calcula, así que
		\[
			\boxed{\;\mathcal{T}_3 \text{ es el único no causal.}\;}
		\]
		Vale registrar un caso particular que se repite en todos los ejercicios: un sistema
		sin memoria es automáticamente causal, porque usar solo $x[n]$ ya cumple la
		condición $k\le n$.

		\item Los sistemas $\mathcal{T}_1$ y
		$\mathcal{T}_3$ combinan muestras desplazadas de la entrada con coeficientes
		constantes, que es la forma que la regla práctica reconoce, así que son LIT. De
		todos modos conviene verificarlo con la definición al menos una vez.

		\emph{Linealidad de $\mathcal{T}_1$.} Se aplica la regla a la entrada combinada y se
		reagrupan los términos de cada señal:
		\[
			\begin{aligned}
			\mathcal{T}_1\{\alpha x_1[n]+\beta x_2[n]\}
			&= \bigl(\alpha x_1[n]+\beta x_2[n]\bigr) - 2\bigl(\alpha x_1[n-2]+\beta x_2[n-2]\bigr) \\
			&= \alpha\bigl(x_1[n]-2x_1[n-2]\bigr) + \beta\bigl(x_2[n]-2x_2[n-2]\bigr) \\
			&= \alpha\,\mathcal{T}_1\{x_1[n]\} + \beta\,\mathcal{T}_1\{x_2[n]\}. \quad\checkmark
			\end{aligned}
		\]

		\emph{Invarianza de $\mathcal{T}_1$.} Se calculan por separado las dos secuencias que
		la definición compara. Llamando $x_d[n]=x[n-n_0]$ a la entrada desplazada, cuya
		muestra dos lugares atrás es $x_d[n-2]=x[n-2-n_0]$,
		\[
			\begin{aligned}
			\mathcal{T}_1\{x_d[n]\} &= x_d[n]-2x_d[n-2] = x[n-n_0]-2\,x[n-n_0-2], \\
			y[n-n_0] &= x[n-n_0]-2\,x[(n-n_0)-2] = x[n-n_0]-2\,x[n-n_0-2],
			\end{aligned}
		\]
		y las dos expresiones coinciden.

		\emph{$\mathcal{T}_3$.} La linealidad se sigue de que la suma se reparte en los dos
		términos y cada escalar sale como factor común de su propia suma. Para la invarianza,
		la salida ante la entrada desplazada es
		\[
			\mathcal{T}_3\{x_d[n]\} = \sum_{k=n-2}^{n+2} x[k-n_0],
		\]
		y el cambio de variable $m=k-n_0$ traslada los dos extremos de la suma, que pasan a
		ser $m=n-n_0-2$ y $m=n-n_0+2$. Lo que queda es exactamente $y[n-n_0]$.

		\emph{$\mathcal{T}_2$ no es lineal.} La regla multiplica dos muestras de la entrada,
		así que al escalar la entrada por $\alpha$ los dos factores se escalan a la vez,
		\[
			\mathcal{T}_2\{\alpha\,x[n]\} = \bigl(\alpha x[n]\bigr)\bigl(\alpha x[n-1]\bigr)
			= \alpha^2\,y[n],
		\]
		y $\alpha^2 \neq \alpha$ salvo para $\alpha=0$ y $\alpha=1$. Con un caso numérico se
		ve mejor. Tomando $x[n]=\delta[n]+\delta[n-1]$, el único producto no nulo ocurre en
		$n=1$, donde se cruzan las dos muestras unitarias, de modo que $y[n]=\delta[n-1]$.
		Duplicando la entrada, cada factor se duplica y la salida se cuadruplica,
		\[
			\mathcal{T}_2\{2x[n]\} = 4\,\delta[n-1]
			\quad\text{frente a}\quad
			2\,\mathcal{T}_2\{x[n]\} = 2\,\delta[n-1].
		\]
		La homogeneidad falla y con ella la linealidad.

		\emph{$\mathcal{T}_2$ sí es invariante.} Su regla no menciona al índice $n$ más que a
		través de las muestras que toma,
		\[
			\mathcal{T}_2\{x_d[n]\} = x_d[n]\,x_d[n-1] = x[n-n_0]\,x[n-1-n_0] = y[n-n_0].
		\]
		Un sistema puede entonces ser invariante sin ser lineal, y este es el ejemplo.

		\emph{$\mathcal{T}_4$ es lineal.} El factor $(n+1)$ no depende de la entrada, así que
		multiplica por igual a los dos términos de la combinación y los escalares salen
		afuera,
		\[
			(n+1)\bigl(\alpha x_1[n]+\beta x_2[n]\bigr)
			= \alpha\,(n+1)x_1[n] + \beta\,(n+1)x_2[n]. \quad\checkmark
		\]

		\emph{$\mathcal{T}_4$ no es invariante.} Aquí está la trampa del ejercicio, porque el
		factor variable no rompe la linealidad pero sí la invarianza. Ante la entrada
		desplazada el sistema sigue multiplicando por $(n+1)$, mientras que la salida
		desplazada lleva el factor evaluado en el índice desplazado,
		\[
			\begin{aligned}
			\mathcal{T}_4\{x[n-n_0]\} &= (n+1)\,x[n-n_0], \\
			y[n-n_0] &= (n-n_0+1)\,x[n-n_0].
			\end{aligned}
		\]
		Los dos factores difieren en $n_0$. Para exhibirlo sobre valores concretos, tomemos
		$x[n]=\delta[n]$ y $n_0=1$. La salida ante el impulso es
		$y[n]=(n+1)\delta[n]=\delta[n]$, porque el impulso solo es no nulo en $n=0$, donde el
		factor vale $1$; desplazándola, $y[n-1]=\delta[n-1]$. En cambio, la salida ante el
		impulso desplazado es
		\[
			\mathcal{T}_4\{\delta[n-1]\} = (n+1)\,\delta[n-1] = 2\,\delta[n-1],
		\]
		porque ahora la muestra no nula cae en $n=1$ y allí el factor vale $2$. Las dos
		secuencias difieren, así que el sistema no es invariante en el tiempo.

		\item Se parte de una cota $|x[n]| \le M_x$ para toda
		muestra y se busca una cota para la salida, acotando cada término con la desigualdad
		triangular:
		\[
			\begin{aligned}
			\mathcal{T}_1&: & |y[n]| &\le |x[n]| + 2\,|x[n-2]| \le 3M_x, \\
			\mathcal{T}_2&: & |y[n]| &\le |x[n]|\,|x[n-1]| \le M_x^2, \\
			\mathcal{T}_3&: & |y[n]| &\le \sum_{k=n-2}^{n+2} |x[k]| \le 5M_x.
			\end{aligned}
		\]
		Los tres tienen una cota finita construida solo con $M_x$, y por lo tanto los tres
		son estables. La cota de $\mathcal{T}_2$ crece con el cuadrado y la de
		$\mathcal{T}_3$ con la cantidad de términos de la suma, pero ninguna de las dos
		depende de $n$, que es lo que la estabilidad BIBO exige.

		El sistema $\mathcal{T}_4$, en cambio, no es estable, y para probarlo alcanza con
		exhibir una sola entrada acotada cuya salida no lo esté. Con la entrada constante
		$x[n]=1$, acotada por $M_x=1$, la salida vale
		\[
			y[n] = (n+1)\cdot 1 = n+1,
		\]
		que supera cualquier cota que se fije de antemano en cuanto $n$ es suficientemente
		grande. El motivo es el mismo factor que rompía la invarianza, porque un coeficiente
		que crece con el índice termina amplificando sin límite.

		\item El único que no puede usarse en tiempo
		real es $\mathcal{T}_3$, y es el único no causal, lo cual no es casualidad. En un
		procesamiento en tiempo real las muestras ingresan de a una al ritmo del muestreo, y
		en el instante $n$ el sistema dispone de $x[n]$, de las anteriores y de las salidas
		ya calculadas. Para producir $y[n]$, $\mathcal{T}_3$ necesita $x[n+1]$ y $x[n+2]$,
		dos muestras que en ese momento no existen todavía.

		La regla sí es utilizable cuando la secuencia está completa en memoria antes de
		empezar ---un archivo de audio, una imagen, el registro de una campaña de
		medición---, porque en esa situación las muestras posteriores están disponibles y el
		sistema puede leerlas. Eso es procesamiento fuera de línea, y ahí un sistema no
		causal es una opción legítima. Los otros tres son causales y funcionan en los dos
		modos.
	\end{enumerate}

	La Tabla~\ref{tab:sol_ej1} reúne las cuatro clasificaciones.

	\begin{table}[H]
		\centering
		\caption{Clasificación de los cuatro sistemas del ejercicio 1.}
		\label{tab:sol_ej1}
		\begin{tabular}{lccccc}
			\toprule
			Sistema & Memoria & Causal & Lineal & Invariante & Estable \\
			\midrule
			$\mathcal{T}_1: x[n]-2x[n-2]$          & sí & sí & sí & sí & sí \\
			$\mathcal{T}_2: x[n]\,x[n-1]$          & sí & sí & no & sí & sí \\
			$\mathcal{T}_3: \sum_{k=n-2}^{n+2}x[k]$ & sí & no & sí & sí & sí \\
			$\mathcal{T}_4: (n+1)\,x[n]$            & no & sí & sí & no & no \\
			\bottomrule
		\end{tabular}
	\end{table}

	Solo $\mathcal{T}_1$ y $\mathcal{T}_3$ son LIT, y son justamente los dos que la
	regla práctica reconocía a simple vista. Los otros dos muestran que las propiedades son
	independientes entre sí: $\mathcal{T}_2$ es invariante sin ser lineal y $\mathcal{T}_4$
	es lineal sin ser invariante.

	% ============================================================================
	\tema{Respuesta al impulso y convolución}
	% ============================================================================

	La \textbf{respuesta al impulso} de un sistema es la salida que produce ante la entrada
	$\delta[n]$, con el sistema en reposo,
	\[
		h[n] = \mathcal{T}\{\delta[n]\},
	\]
	y se obtiene corriendo la regla del sistema sobre esa entrada. Cuando el sistema es
	LIT, $h[n]$ lo caracteriza por completo, porque la salida ante cualquier entrada es la
	\textbf{suma de convolución}
	\[
		y[n] = \sum_{k=-\infty}^{\infty} x[k]\,h[n-k] \;=\; x[n]*h[n].
	\]

	Esa suma admite dos lecturas, y las dos llegan al mismo número.

	\begin{itemize}[nosep]
		\item \textbf{Copias desplazadas.} Fijando $k$, el sumando $x[k]\,h[n-k]$ es una
		      copia de $h[n]$ escalada por $x[k]$ y desplazada hasta la posición $k$. La
		      salida es la superposición de una copia por cada muestra no nula de la
		      entrada, y se tabula con una fila por copia y una columna por instante.
		\item \textbf{Invertir y desplazar.} Para cada $n$, se grafican $x[k]$ y $h[n-k]$
		      sobre el eje $k$, se multiplica muestra a muestra y se suman los productos.
		      La secuencia $h[n-k]$ se construye invirtiendo $h$ y desplazándola $n$
		      lugares, y funciona como una ventana de pesos que pondera las muestras de la entrada.
	\end{itemize}

	Dos consecuencias de la tabla de copias se usan todo el tiempo. La primera es que la
	convolución equivale al \textbf{producto de polinomios} cuyos coeficientes son las
	muestras de cada secuencia. La segunda es el \textbf{largo del soporte}: si $x[n]$
	ocupa $L_x$ muestras y $h[n]$ ocupa $L_h$, la salida ocupa $L_x+L_h-1$.

	% ----------------------------------------------------------------------------
	\ejercicio{2: convolución de dos secuencias finitas por tres caminos}

	\begin{consigna}
	Un sistema LIT responde a la ecuación
	\[
		y[n] = 2\,x[n] + x[n-1] - x[n-3],
	\]
	y se lo excita con la entrada
	\[
		x[n] = \{\, \underset{\uparrow}{1} \quad 2 \quad -1 \,\}.
	\]
	\begin{enumerate}
		\item Obtener y graficar $h[n]$, indicando su soporte.
		\item Calcular $y[n]$ armando la tabla de copias desplazadas de
		$h[n]$, con una fila por muestra no nula de la entrada.
		\item Recalcular la salida completa con el procedimiento de invertir y
		desplazar, graficando $x[k]$ y $h[n-k]$ en cada instante, y verificar que
		coincide con la última fila de la tabla.
		\item Contrastar el largo del soporte de la salida con el que
		predice la regla $L_x + L_h - 1$.
		\item Rehacer el cálculo como producto de polinomios y comparar
		los coeficientes con la última fila de la tabla.
	\end{enumerate}
	\end{consigna}

	La entrada tiene la flecha en la primera muestra, de modo que $x[0]=1$, $x[1]=2$,
	$x[2]=-1$ y vale cero fuera del soporte $[0,2]$.

	\begin{enumerate}
		\item Se obtiene reemplazando la entrada por
		$\delta[n]$ en la ecuación del sistema. Cada término $x[n-k]$ pasa a ser
		$\delta[n-k]$, así que
		\[
			h[n] = 2\,\delta[n] + \delta[n-1] - \delta[n-3].
		\]
		Cada impulso aporta un valor en su propio índice, con lo que la respuesta vale $2$
		en $n=0$, $1$ en $n=1$, $0$ en $n=2$ y $-1$ en $n=3$,
		\[
			\boxed{\;h[n] = \{\, \underset{\uparrow}{2} \quad 1 \quad 0 \quad -1 \,\},
			\qquad \text{soporte } [0,3],\quad L_h = 4.\;}
		\]
		El valor nulo en $n=2$ está dentro del soporte, porque el soporte es el intervalo de
		índices fuera del cual la secuencia se anula y no el conjunto de sus muestras no
		nulas. Los paneles (a) y (b) de la Figura~\ref{fig:sol_ej2} muestran la entrada y la
		respuesta al impulso.

		\item La entrada tiene tres muestras no
		nulas, así que la salida es la suma de tres copias de $h[n]$. Cada copia se escala
		por el valor de la muestra que la produce y arranca en la posición de esa muestra.
		La Tabla~\ref{tab:sol_ej2} ordena la cuenta con una fila por copia y una columna por
		instante.

		\begin{table}[H]
			\centering
			\caption{Ejercicio~2. Cada fila es la copia de $h[n]$ que aporta una muestra de
			la entrada, escalada por su valor y desplazada hasta su posición. Las casillas
			vacías son ceros y la última fila suma cada columna.}
			\label{tab:sol_ej2}
			\renewcommand{\arraystretch}{1.3}
			\begin{tabular}{lcccccc}
				\toprule
				Aporte & $n=0$ & $n=1$ & $n=2$ & $n=3$ & $n=4$ & $n=5$ \\
				\midrule
				$x[0]\,h[n]   = \phantom{-}1\cdot h[n]$   & $2$ & $1$ & $0$  & $-1$ &      &     \\
				$x[1]\,h[n-1] = \phantom{-}2\cdot h[n-1]$ &     & $4$ & $2$  & $0$  & $-2$ &     \\
				$x[2]\,h[n-2] = -1\cdot h[n-2]$           &     &     & $-2$ & $-1$ & $0$  & $1$ \\
				\midrule
				$y[n]$ & $2$ & $5$ & $0$ & $-2$ & $-2$ & $1$ \\
				\bottomrule
			\end{tabular}
		\end{table}

		La primera fila es $h[n]$ sin modificar, porque $x[0]=1$ no la escala y no hay
		desplazamiento. La segunda repite los mismos cuatro valores duplicados y desplazados
		un lugar, y la tercera los repite cambiados de signo y desplazados dos lugares.
		Sumando cada columna,
		\[
			\boxed{\;y[n] = \{\, \underset{\uparrow}{2} \quad 5 \quad 0 \quad -2 \quad -2 \quad 1 \,\}.\;}
		\]
		El valor nulo en $n=2$ tiene otro origen que el de $h[n]$, porque aquí los tres
		aportes de esa columna sí son no nulos y se cancelan entre sí, $0+2-2=0$. El panel
		(c) de la Figura~\ref{fig:sol_ej2} muestra la salida.

		\item La segunda lectura calcula un valor
		de la salida por vez. Se elige el instante $n$, se arma el peso $h[n-k]$ que le
		corresponde a cada muestra de la entrada y se suman los productos,
		\[
			y[n] = \sum_{k=-\infty}^{\infty} x[k]\,h[n-k].
		\]
		Recorriendo así un instante tras otro se obtiene la salida completa, que es lo que
		pide el inciso.

		Lo único que cambia de un instante al siguiente es la posición de $h[n-k]$ frente a
		la entrada, así que el barrido se sigue mejor sobre la gráfica que sobre la fórmula.
		La Figura~\ref{fig:sol_ej2c} arma las tres secuencias para los seis instantes en que
		la salida es no nula, uno por columna.

		\begin{figure}[H]
			\centering
			\includestandalone[width=\linewidth]{figures/fig_sol_ej2c/fig_sol_ej2c}
			\caption{Ejercicio~2c. Un instante por columna, con la entrada $x[k]$ arriba, los
			pesos $h[n-k]$ en el medio y la salida $y[n]$ abajo. En cada columna la muestra
			de salida que se calcula va en tono pleno y las ya calculadas atenuadas.}
			\label{fig:sol_ej2c}
		\end{figure}

		La fila de arriba es la entrada $x[k]$, la misma en las seis columnas. En la fila del
		medio está $h[n-k]$, la respuesta al impulso reflejada y desplazada según el valor de
		$n$, de modo que avanza un lugar de una columna a la siguiente. La franja gris cubre
		los cuatro índices donde los pesos son no nulos, y fuera de ella todos los productos
		valen cero, así que en cada instante solo participan las muestras de la entrada que
		quedan bajo la franja. La fila de abajo es la salida, que gana una muestra por
		columna.

		En cada columna se multiplica cada muestra de la entrada por el peso que le cae
		encima y se suman los productos,
		\begin{align*}
			y[0] &= x[0]\,h[0] = 1\cdot 2 = 2, \\
			y[1] &= x[0]\,h[1] + x[1]\,h[0] = 1 + 4 = 5, \\
			y[2] &= x[0]\,h[2] + x[1]\,h[1] + x[2]\,h[0] = 0 + 2 - 2 = 0, \\
			y[3] &= x[0]\,h[3] + x[1]\,h[2] + x[2]\,h[1] = -1 + 0 - 1 = -2, \\
			y[4] &= x[1]\,h[3] + x[2]\,h[2] = -2 + 0 = -2, \\
			y[5] &= x[2]\,h[3] = (-1)\cdot(-1) = 1.
		\end{align*}
		El cero de $n=2$ sale de que los dos últimos aportes se cancelan entre sí, aunque
		ninguna de las tres muestras de la entrada sea nula. En $n=6$ la ventana ya pasó de
		largo el soporte de la entrada, todos los productos valen cero y la salida se apaga.
		Los seis valores obtenidos son
		\[
			\boxed{\;y[n] = \{\, \underset{\uparrow}{2} \quad 5 \quad 0 \quad -2 \quad -2 \quad 1 \,\},\;}
		\]
		la última fila de la Tabla~\ref{tab:sol_ej2}. Las dos lecturas recorren los mismos
		productos en distinto orden, la tabla acumulando las copias que caen en cada instante
		y la ventana ponderando las muestras de la entrada que lo alimentan.

		\item La entrada ocupa $L_x=3$ muestras y la respuesta
		al impulso $L_h=4$, de modo que la regla predice
		\[
			L_y = L_x + L_h - 1 = 3 + 4 - 1 = 6 \ \text{muestras}.
		\]
		La salida obtenida ocupa el soporte $[0,5]$, que son exactamente seis muestras. El
		conteo se lee sobre la tabla: las tres primeras casillas ocupadas caen en columnas
		distintas y cubren $L_x=3$ columnas, y la última fila se prolonga $L_h-1=3$ columnas
		más hacia la derecha.

		\item A cada secuencia causal se le asocia el
		polinomio en el marcador $D$ cuyos coeficientes son sus muestras, con $D^k$ ocupando
		el lugar del índice $k$:
		\[
			X(D) = 1 + 2D - D^2, \qquad H(D) = 2 + D - D^3.
		\]
		Multiplicando cada término del primer factor por cada término del segundo y
		reuniendo después los que comparten exponente,
		\[
			\begin{aligned}
			Y(D) &= (1+2D-D^2)(2+D-D^3) \\
			     &= 2 + D - D^3 \;+\; 4D + 2D^2 - 2D^4 \;-\; 2D^2 - D^3 + D^5 \\
			     &= 2 + 5D + 0\,D^2 - 2D^3 - 2D^4 + D^5.
			\end{aligned}
		\]
		Los seis coeficientes son la última fila de la Tabla~\ref{tab:sol_ej2}, en el mismo
		orden. El marcador $D$ lleva por sí solo la contabilidad de qué muestra de $x$ se
		cruza con qué muestra de $h$, que es la misma que en la tabla llevaba el
		desplazamiento progresivo de cada fila.
	\end{enumerate}

	\begin{figure}[H]
		\centering
		\includestandalone[width=0.45\linewidth]{figures/fig_sol_ej2/fig_sol_ej2}
		\caption{Ejercicio~2. (a) La entrada $x[n]$, soporte $[0,2]$. (b) La respuesta al
		impulso $h[n]$, soporte $[0,3]$. (c) La salida $y[n]$, soporte $[0,5]$. Los tres
		paneles comparten la escala vertical.}
		\label{fig:sol_ej2}
	\end{figure}

	% ----------------------------------------------------------------------------
	\ejercicio{3: convoluciones con escalones, impulsos y exponenciales}

	\begin{consigna}
	Sean las convoluciones
	\begin{align*}
		y_1[n] &= \bigl(u[n]-u[n-4]\bigr) * \bigl(u[n]-u[n-3]\bigr), \\
		y_2[n] &= \left(\tfrac{1}{2}\right)^{\!n} u[n] *
		          \left(\delta[n]-\tfrac{1}{2}\,\delta[n-1]\right), \\
		y_3[n] &= \left(\tfrac{1}{2}\right)^{\!n} u[n] *
		          \left(\tfrac{1}{4}\right)^{\!n} u[n].
	\end{align*}
	\begin{enumerate}
		\item Calcular y graficar $y_1[n]$, indicando su soporte.
		\item Calcular $y_2[n]$ e indicar qué relación guardan entre sí
		los dos sistemas cuyas respuestas al impulso se convolucionaron.
		\item Calcular $y_3[n]$ en forma cerrada e indicar cuál de las dos
		exponenciales gobierna el decaimiento cuando $n$ es grande.
	\end{enumerate}
	\end{consigna}

	\begin{enumerate}
		\item La diferencia de dos escalones desplazados
		es un pulso, y lo primero es escribir las dos secuencias por extensión. El escalón
		$u[n]$ vale uno desde $n=0$ en adelante y $u[n-4]$ vale uno desde $n=4$, así que la
		resta deja cuatro unos,
		\[
			p_1[n] = u[n]-u[n-4] = \{\, \underset{\uparrow}{1} \quad 1 \quad 1 \quad 1 \,\},
			\qquad
			p_2[n] = u[n]-u[n-3] = \{\, \underset{\uparrow}{1} \quad 1 \quad 1 \,\},
		\]
		de soportes $[0,3]$ y $[0,2]$, con $L_1=4$ y $L_2=3$. La salida ocupa entonces
		$4+3-1=6$ muestras, sobre el soporte $[0,5]$.

		El pulso corto tiene tres muestras no nulas, así que $y_1[n]$ es la suma de tres
		copias del pulso largo, cada una escalada por la muestra que la produce y arrancada
		en la posición de esa muestra. La Tabla~\ref{tab:sol_ej3} las ordena con una fila por
		copia y una columna por instante.

		\begin{table}[H]
			\centering
			\caption{Ejercicio~3a. Las tres copias de $p_1[n]$ que aporta cada muestra de
			$p_2[n]$. Como las tres muestras valen uno, ninguna copia se escala y las filas
			son el mismo bloque de cuatro unos corrido un lugar por vez. Las casillas vacías
			son ceros y la última fila suma cada columna.}
			\label{tab:sol_ej3}
			\renewcommand{\arraystretch}{1.3}
			\begin{tabular}{lcccccc}
				\toprule
				Aporte & $n=0$ & $n=1$ & $n=2$ & $n=3$ & $n=4$ & $n=5$ \\
				\midrule
				$p_2[0]\,p_1[n]   = p_1[n]$   & $1$ & $1$ & $1$ & $1$ &     &     \\
				$p_2[1]\,p_1[n-1] = p_1[n-1]$ &     & $1$ & $1$ & $1$ & $1$ &     \\
				$p_2[2]\,p_1[n-2] = p_1[n-2]$ &     &     & $1$ & $1$ & $1$ & $1$ \\
				\midrule
				$y_1[n]$ & $1$ & $2$ & $3$ & $3$ & $2$ & $1$ \\
				\bottomrule
			\end{tabular}
		\end{table}

		Sumando cada columna,
		\[
			\boxed{\;y_1[n] = \{\, \underset{\uparrow}{1} \quad 2 \quad 3 \quad 3 \quad 2 \quad 1 \,\},
			\qquad \text{soporte } [0,5].\;}
		\]
		El resultado es la forma trapezoidal del panel (a) de la
		Figura~\ref{fig:sol_ej3}, con una rampa de subida mientras el solapamiento crece, una
		meseta mientras el pulso corto está completamente cubierto y una rampa de bajada
		simétrica. Es la versión discreta del trapecio que producen dos pulsos rectangulares
		en tiempo continuo.

		\item Convolucionar con un impulso
		desplazado es desplazar, porque $x[n]*\delta[n-n_0]=x[n-n_0]$. La segunda secuencia
		tiene dos impulsos, así que por distributividad la convolución se separa en dos
		términos:
		\[
			y_2[n] = \left(\tfrac{1}{2}\right)^{\!n} u[n] * \delta[n]
			       \;-\; \tfrac{1}{2}\left(\tfrac{1}{2}\right)^{\!n} u[n] * \delta[n-1].
		\]
		El primer término queda intacto, porque el impulso es el elemento neutro. El segundo
		es la misma secuencia desplazada un lugar y multiplicada por $1/2$,
		\[
			\tfrac{1}{2}\left(\tfrac{1}{2}\right)^{\!n-1} u[n-1]
			= \left(\tfrac{1}{2}\right)^{\!n} u[n-1],
		\]
		donde el factor $1/2$ de afuera se absorbió en el exponente, que pasa de $n-1$ a $n$.
		Los dos términos comparten entonces el mismo factor exponencial y difieren solo en el
		escalón,
		\[
			y_2[n] = \left(\tfrac{1}{2}\right)^{\!n}\bigl(u[n]-u[n-1]\bigr)
			       = \left(\tfrac{1}{2}\right)^{\!n} \delta[n],
		\]
		y como el impulso solo es no nulo en $n=0$, donde la exponencial vale uno,
		\[
			\boxed{\;y_2[n] = \delta[n].\;}
		\]
		La convolución de las dos respuestas al impulso es la respuesta al impulso de la
		cascada de los dos sistemas, y aquí esa cascada devuelve el impulso sin tocarlo. Los
		dos sistemas son entonces \textbf{mutuamente inversos}: encadenados en cualquier
		orden reconstruyen la entrada, así que cada uno deshace exactamente el efecto del otro.

		\item Los dos escalones recortan el rango del
		índice. El primer factor anula los términos con $k<0$ y el segundo los que tienen
		$n-k<0$, es decir $k>n$, de modo que sobreviven los índices $0 \le k \le n$ y la suma
		queda finita:
		\[
			y_3[n] = \sum_{k=0}^{n} \left(\tfrac{1}{2}\right)^{\!k}
			         \left(\tfrac{1}{4}\right)^{\!n-k}.
		\]
		Contar términos ya no alcanza, porque cada uno vale algo distinto. Lo que sí se puede
		es separar el exponente del segundo factor, $\left(\tfrac14\right)^{n-k}
		=\left(\tfrac14\right)^{n}\left(\tfrac14\right)^{-k}$, y sacar afuera de la suma lo
		que no depende de $k$:
		\[
			y_3[n] = \left(\tfrac{1}{4}\right)^{\!n} \sum_{k=0}^{n} \left(\frac{1/2}{1/4}\right)^{\!k}
			       = \left(\tfrac{1}{4}\right)^{\!n} \sum_{k=0}^{n} 2^{k}.
		\]
		Adentro quedó una suma geométrica de razón $2$ y $n+1$ términos, cuya forma cerrada
		es $(2^{\,n+1}-1)/(2-1) = 2^{\,n+1}-1$. Reemplazando y repartiendo el factor de
		afuera,
		\[
			y_3[n] = \frac{2^{\,n+1}-1}{4^{\,n}}
			       = 2\,\frac{2^{\,n}}{4^{\,n}} - \frac{1}{4^{\,n}},
		\]
		y como $2^n/4^n = (1/2)^n$, la salida es
		\[
			\boxed{\;y_3[n] = \left[\,2\left(\tfrac{1}{2}\right)^{\!n}
			         - \left(\tfrac{1}{4}\right)^{\!n}\right] u[n].\;}
		\]
		Los dos primeros valores comprueban la fórmula contra la suma original. Para $n=0$ hay
		un solo término, $\left(\tfrac12\right)^0\left(\tfrac14\right)^0=1$, y la fórmula
		vale $2-1=1$. Para $n=1$ la suma tiene dos términos,
		$1\cdot\tfrac14 + \tfrac12\cdot 1 = \tfrac34$, y la fórmula arroja
		$2\cdot\tfrac12-\tfrac14=\tfrac34$.

		El resultado no es una exponencial sino una combinación de las dos, y ninguna de las
		dos bases desapareció. Como $1/4 < 1/2$, el término $\left(\tfrac14\right)^n$ se
		achica mucho más rápido y se vuelve despreciable frente al otro, de modo que para $n$
		grande
		\[
			y_3[n] \approx 2\left(\tfrac{1}{2}\right)^{\!n}.
		\]
		El decaimiento lo gobierna entonces la exponencial \emph{más lenta} de las dos, la de
		base $1/2$. El panel (b) de la Figura~\ref{fig:sol_ej3} lo muestra sobre la gráfica:
		la curva punteada es $2\,(1/2)^n$, y las muestras de $y_3[n]$ se le acercan a partir
		de $n=3$. La lectura de sistemas es que al encadenar dos sistemas de respuestas
		$(1/2)^n u[n]$ y $(1/4)^n u[n]$, quien fija la velocidad con que se apaga la cascada
		es la etapa más lenta.
	\end{enumerate}

	\begin{figure}[H]
		\centering
		\includestandalone[width=0.62\linewidth]{figures/fig_sol_ej3/fig_sol_ej3}
		\caption{Ejercicio~3. (a) $y_1[n]$, el trapecio que producen dos pulsos
		rectangulares, soporte $[0,5]$. (b) $y_3[n]$, junto con la exponencial
		$2\,(1/2)^n$ punteada, que gobierna el decaimiento para $n$ grande.}
		\label{fig:sol_ej3}
	\end{figure}

	% ----------------------------------------------------------------------------
	\ejercicio{4: la respuesta al escalón}

	\begin{consigna}
	Un sistema LIT tiene respuesta al impulso
	$h[n]=\left(\tfrac{1}{2}\right)^{n} u[n]$. Se llama $g[n]$ a la salida
	que produce ante la entrada $u[n]$.
	\begin{enumerate}
		\item Resolver la suma de convolución y escribir $g[n]$ en forma
		cerrada.
		\item Mostrar, sobre esa misma suma, que $g[n]$ es la suma
		acumulada de los valores $h[k]$ con $k \le n$.
		\item Verificar que $h[n]=g[n]-g[n-1]$ y explicar por qué es la
		primera diferencia la que recupera $h[n]$.
		\item Determinar hacia qué valor tiende $g[n]$ al aumentar $n$ y
		relacionar ese límite con la estabilidad del sistema.
	\end{enumerate}
	\end{consigna}

	\begin{enumerate}
		\item La salida ante el escalón es la convolución entre
		$h[n]$ y $u[n]$,
		\[
			g[n] = \sum_{k=-\infty}^{\infty} h[k]\,u[n-k]
			     = \sum_{k=-\infty}^{\infty} \left(\tfrac{1}{2}\right)^{\!k} u[k]\,u[n-k].
		\]
		El escalón $u[k]$ anula los términos con $k<0$ y el escalón $u[n-k]$ anula aquellos
		en que $n-k<0$, es decir $k>n$. Sobreviven los índices $0\le k\le n$, siempre que
		$n\ge 0$; si $n<0$ el rango queda vacío y la salida es nula. Para $n \ge 0$,
		\[
			g[n] = \sum_{k=0}^{n} \left(\tfrac{1}{2}\right)^{\!k}
			     = \frac{1-\left(\tfrac{1}{2}\right)^{n+1}}{1-\tfrac{1}{2}}
			     = 2\left[\,1-\left(\tfrac{1}{2}\right)^{\!n+1}\right],
		\]
		donde se aplicó la forma cerrada de la suma geométrica finita de razón $1/2$ y $n+1$
		términos. Distribuyendo el $2$, el factor $\left(\tfrac12\right)^{n+1}$ pierde una
		potencia y queda
		\[
			\boxed{\;g[n] = \left[\,2-\left(\tfrac{1}{2}\right)^{\!n}\right] u[n].\;}
		\]
		Los primeros valores son $g[0]=1$, $g[1]=1{,}5$, $g[2]=1{,}75$ y $g[3]=1{,}875$, que
		es la sucesión creciente del panel (a) de la Figura~\ref{fig:sol_ej4}.

		\item La respuesta al escalón es la suma acumulada de $h[n]$, y el paso está a la
		vista en la cuenta anterior. Vale para cualquier sistema LIT, no solo para este. Volviendo a la
		suma antes de reemplazar la exponencial,
		\[
			g[n] = \sum_{k=-\infty}^{\infty} h[k]\,u[n-k],
		\]
		el factor $u[n-k]$ vale uno cuando $k \le n$ y cero cuando $k>n$. No escala ningún
		término, solo decide cuáles participan, de modo que la suma se reduce a
		\[
			\boxed{\;g[n] = \sum_{k=-\infty}^{n} h[k].\;}
		\]
		La respuesta al escalón en cada instante es entonces el total acumulado de la
		respuesta al impulso hasta ese instante. Sobre este ejemplo se verifica de inmediato,
		porque $g[2]=h[0]+h[1]+h[2]=1+0{,}5+0{,}25=1{,}75$, el mismo valor de la forma
		cerrada.

		\item Con $g[n]$ escrita como suma
		acumulada, la resta de dos totales consecutivos cancela todo lo acumulado hasta el
		instante anterior y deja un único término,
		\[
			g[n]-g[n-1] = \sum_{k=-\infty}^{n} h[k] \;-\; \sum_{k=-\infty}^{n-1} h[k] = h[n].
		\]
		Sobre la forma cerrada se comprueba con números. Para $n \ge 1$,
		\[
			g[n]-g[n-1]
			= \left[2-\left(\tfrac{1}{2}\right)^{\!n}\right]
			- \left[2-\left(\tfrac{1}{2}\right)^{\!n-1}\right]
			= \left(\tfrac{1}{2}\right)^{\!n-1} - \left(\tfrac{1}{2}\right)^{\!n}
			= \left(\tfrac{1}{2}\right)^{\!n},
		\]
		donde el último paso sale de que $\left(\tfrac12\right)^{n-1}=2\left(\tfrac12\right)^{n}$,
		así que la resta deja una sola copia. En $n=0$ la cuenta es $g[0]-g[-1]=1-0=1=h[0]$,
		y para $n<0$ ambos términos son nulos. Los valores recuperados son los del panel (b)
		de la Figura~\ref{fig:sol_ej4}.

		Que sea la primera diferencia y no otra operación tiene una razón que no depende de
		este sistema. El impulso es la primera diferencia del escalón,
		$\delta[n]=u[n]-u[n-1]$, y el sistema es LIT. Por invarianza temporal, la entrada
		$u[n-1]$ produce la salida $g[n-1]$; por linealidad, la entrada $u[n]-u[n-1]$ produce
		la salida $g[n]-g[n-1]$. Esa entrada es $\delta[n]$ y esa salida es, por definición,
		$h[n]$. La operación que en la entrada convierte el escalón en impulso es la misma
		que en la salida convierte $g[n]$ en $h[n]$.

		\item Al aumentar $n$, el término
		$\left(\tfrac12\right)^n$ se achica hacia cero y
		\[
			\boxed{\;\lim_{n\to\infty} g[n] = 2.\;}
		\]
		El valor no es casual, porque el inciso (b) mostró que $g[n]$ acumula los valores de
		$h[k]$, de modo que su límite es la suma de todos ellos,
		\[
			\sum_{n=-\infty}^{\infty} h[n] = \sum_{n=0}^{\infty} \left(\tfrac{1}{2}\right)^{\!n}
			= \frac{1}{1-\tfrac{1}{2}} = 2.
		\]
		El criterio de estabilidad BIBO pide que $h[n]$ sea absolutamente sumable, y aquí
		todos los valores son positivos, así que la suma de los módulos es la misma serie y
		también vale $2$. El sistema es estable, y la respuesta al escalón lo exhibe, porque
		ante una entrada acotada por $1$ la salida se estaciona en un valor finito en lugar de
		crecer sin límite. Si en cambio la base fuera mayor o igual que uno, la
		serie no convergería y $g[n]$ crecería indefinidamente ante la misma entrada acotada.
	\end{enumerate}

	\begin{figure}[H]
		\centering
		\includestandalone[width=0.5\linewidth]{figures/fig_sol_ej4/fig_sol_ej4}
		\caption{Ejercicio~4. (a) La respuesta al escalón $g[n]=2-(1/2)^n$, que se acerca a
		la línea punteada de valor $2$. (b) La respuesta al impulso $h[n]=(1/2)^n u[n]$
		recuperada como primera diferencia de $g[n]$.}
		\label{fig:sol_ej4}
	\end{figure}

	% ============================================================================
	\tema{Interconexión de sistemas LIT}
	% ============================================================================

	Las propiedades algebraicas de la convolución se leen directamente sobre las
	interconexiones. Dos sistemas LIT en \textbf{cascada} equivalen a uno solo de respuesta
	al impulso $h_1[n]*h_2[n]$, y por la conmutatividad el orden de las etapas es
	intercambiable. Dos sistemas en \textbf{paralelo}, que reciben la misma entrada y suman
	sus salidas, equivalen a uno de respuesta $h_1[n]+h_2[n]$. Cuando hay etapas mezcladas,
	se reduce de a una por vez, empezando por la más interna.

	% ----------------------------------------------------------------------------
	\ejercicio{5: paralelo en cascada}

	\begin{consigna}
	El sistema de la Figura~\ref{fig:tp7_interconexion} se arma con
	tres sistemas LIT de respuestas al impulso
	\begin{align*}
		h_1[n] &= \delta[n]+\delta[n-1], \\
		h_2[n] &= \delta[n]-\delta[n-1], \\
		h_3[n] &= \left(\tfrac{1}{2}\right)^{\!n} u[n].
	\end{align*}

	\begin{center}
		\includestandalone[width=0.5\linewidth]{../figures/fig_tp7_interconexion/fig_tp7_interconexion}
		\captionof{figure}{Interconexión del ejercicio 5.}
		\label{fig:tp7_interconexion}
	\end{center}

	\begin{enumerate}
		\item Hallar la respuesta al impulso equivalente de las dos ramas
		en paralelo e indicar qué operación aplica esa etapa sobre la
		entrada.
		\item Hallar la respuesta al impulso $h[n]$ del sistema completo.
		\item Clasificar el sistema completo como FIR o IIR y determinar
		si es estable.
		\item Indicar si $h[n]$ cambia al intercambiar el orden de las dos
		etapas de la cascada. Justificar.
	\end{enumerate}
	\end{consigna}

	Antes de operar vale identificar qué sistema es cada rama. La respuesta
	$h_1[n]=\delta[n]+\delta[n-1]$ corresponde a $w[n]=x[n]+x[n-1]$, la suma de la muestra
	actual con la anterior, y $h_2[n]=\delta[n]-\delta[n-1]$ corresponde a
	$w[n]=x[n]-x[n-1]$, la primera diferencia. Las dos ramas operan sobre las mismas dos
	muestras y solo difieren en el signo con que interviene la más vieja.

	\begin{enumerate}
		\item Las ramas reciben la misma entrada y sus
		salidas se suman, así que el equivalente es la suma de las respuestas al impulso:
		\[
			h_{12}[n] = h_1[n]+h_2[n]
			= \bigl(\delta[n]+\delta[n-1]\bigr) + \bigl(\delta[n]-\delta[n-1]\bigr)
			= 2\,\delta[n].
		\]
		Los dos términos en $\delta[n-1]$ llegan con signos opuestos y se cancelan, y quedan
		los dos impulsos del origen sumados,
		\[
			\boxed{\;h_{12}[n] = 2\,\delta[n].\;}
		\]
		Una respuesta al impulso concentrada en el origen corresponde a un sistema sin
		memoria, y el factor $2$ es su ganancia. La etapa en paralelo, entonces,
		\textbf{multiplica la entrada por dos} y no conserva nada del pasado, $w[n]=2x[n]$.
		Leído sobre las dos ramas, sumar $x[n]+x[n-1]$ y $x[n]-x[n-1]$ cancela la muestra
		vieja y duplica la actual, que es la misma cuenta escrita sobre las señales en lugar
		de sobre los impulsos.

		\item La etapa en paralelo está en cascada con $h_3$,
		de modo que la respuesta al impulso total es la convolución de las dos:
		\[
			h[n] = h_{12}[n] * h_3[n] = 2\,\delta[n] * \left(\tfrac{1}{2}\right)^{\!n} u[n].
		\]
		El impulso es el elemento neutro de la convolución y el factor $2$ sale afuera por
		linealidad, así que
		\[
			\boxed{\;h[n] = 2\left(\tfrac{1}{2}\right)^{\!n} u[n]
			= \left(\tfrac{1}{2}\right)^{\!n-1} u[n],\;}
		\]
		donde la segunda forma absorbe el $2$ en el exponente. El sistema completo es la
		misma exponencial de la tercera etapa con el doble de amplitud.

		\item El soporte de $h[n]$ es $[0,\infty)$, porque
		la exponencial no se anula en ningún $n\ge 0$, así que el sistema es \textbf{IIR}. Su
		estabilidad se decide con la sumabilidad absoluta,
		\[
			\sum_{n=-\infty}^{\infty}|h[n]| = 2\sum_{n=0}^{\infty}\left(\tfrac{1}{2}\right)^{\!n}
			= 2 \cdot \frac{1}{1-\tfrac{1}{2}} = 4 < \infty,
		\]
		una serie geométrica de razón $1/2 < 1$, que converge. El sistema completo es
		entonces \textbf{estable}. Las dos primeras ramas no podían comprometer el
		resultado, porque tienen soporte finito y todo sistema FIR es estable; quien decide
		es la tercera etapa.

		\item La respuesta al impulso no cambia. La cascada
		de dos sistemas LIT tiene por respuesta la convolución de las dos, y la convolución
		es conmutativa,
		\[
			h_{12}[n] * h_3[n] = h_3[n] * h_{12}[n],
		\]
		de modo que aplicar primero la ganancia y después la exponencial produce el mismo
		sistema que aplicarlas al revés. Sobre este caso se ve sin cuentas, porque escalar
		por dos antes o después de la etapa exponencial conduce al mismo resultado. Vale
		aclarar que la propiedad exige que las dos etapas sean LIT, porque con un bloque no
		lineal, como el cuadrador, el orden sí cambia la salida.
	\end{enumerate}

	% ============================================================================
	\tema{Propiedades leídas en $h[n]$}
	% ============================================================================

	Como $h[n]$ caracteriza por completo a un sistema LIT, las propiedades del catálogo se
	leen directamente sobre ella.

	\begin{itemize}[nosep]
		\item \textbf{Sin memoria} si y solo si $h[n]=K\,\delta[n]$, concentrada en el
		      origen. Cualquier valor no nulo fuera de $n=0$ implica memoria, y la extensión
		      del soporte mide su alcance.
		\item \textbf{Causal} si y solo si $h[n]=0$ para todo $n<0$, es decir si $h[n]$ es
		      una secuencia causal.
		\item \textbf{Estable BIBO} si y solo si $h[n]$ es absolutamente sumable,
		      $\sum_n |h[n]| < \infty$.
		\item \textbf{FIR} si el soporte de $h[n]$ es finito, \textbf{IIR} si es infinito.
		      Todo sistema FIR es estable, porque una suma de finitos términos finitos
		      converge sin excepción.
	\end{itemize}

	% ----------------------------------------------------------------------------
	\ejercicio{6: cinco respuestas al impulso}

	\begin{consigna}
	Sean las respuestas al impulso
	\begin{align*}
		h_1[n] &= 3\,\delta[n], &
		h_2[n] &= u[n]-u[n-5], \\
		h_3[n] &= \left(\tfrac{3}{2}\right)^{\!n} u[n], &
		h_4[n] &= 2^{n}\,u[-n], \\
		h_5[n] &= \cos\!\left(\tfrac{\pi}{4}\,n\right) u[n]. &
	\end{align*}
	\begin{enumerate}
		\item Indicar cuáles corresponden a sistemas sin memoria y cuáles
		a sistemas causales.
		\item Determinar cuáles corresponden a sistemas estables BIBO.
		\item Clasificar cada sistema como FIR o IIR.
		\item Señalar el sistema que resulta estable sin ser causal y
		explicar por qué las dos condiciones son independientes entre sí.
	\end{enumerate}
	\end{consigna}

	El primer paso es identificar el soporte de cada secuencia, porque de ahí salen tres de
	las cuatro respuestas. Las cuatro primeras se leen de la expresión; la quinta merece
	atención porque el coseno cambia de signo.

	\begin{itemize}[nosep]
		\item $h_1[n]$ es un impulso escalado, soporte $\{0\}$.
		\item $h_2[n]$ es el pulso $\{\,\underset{\uparrow}{1}\ \ 1\ \ 1\ \ 1\ \ 1\,\}$,
		      soporte $[0,4]$, porque $u[n]$ arranca en $n=0$ y $u[n-5]$ lo apaga desde
		      $n=5$.
		\item $h_3[n]$ es una exponencial causal creciente, soporte $[0,\infty)$.
		\item $h_4[n]$ vive del otro lado: $u[-n]$ vale uno cuando $-n\ge 0$, es decir para
		      $n\le 0$, y el soporte es $(-\infty,0]$.
		\item $h_5[n]$ es una sinusoide causal de frecuencia digital $\lambda=\pi/4$
		      rad/muestra ($45^\circ$), soporte $[0,\infty)$.
	\end{itemize}

	\begin{enumerate}
		\item Sin memoria hay uno solo, $h_1[n]=3\delta[n]$,
		que tiene la forma $K\delta[n]$ con $K=3$ y corresponde al sistema $y[n]=3x[n]$. Los
		otros cuatro tienen valores no nulos fuera del origen y por lo tanto tienen memoria.

		Para la causalidad hay que mirar el semieje negativo. Las secuencias $h_1$, $h_2$,
		$h_3$ y $h_5$ se anulan en todo $n<0$, así que los cuatro sistemas son causales. La
		excepción es $h_4[n]=2^n u[-n]$, que es no nula justamente para $n\le 0$ y vale, por
		ejemplo, $h_4[-3]=1/8 \neq 0$. Ese sistema responde antes de que el impulso ingrese,
		de modo que \textbf{no es causal}.

		\item Hay que sumar los módulos, uno por uno.

		\emph{$h_1$:} la suma tiene un único término, $\sum_n |h_1[n]| = 3 < \infty$.
		\textbf{Estable.}

		\emph{$h_2$:} cinco unos, $\sum_n |h_2[n]| = 5 < \infty$. \textbf{Estable.}

		\emph{$h_3$:} serie geométrica de razón $3/2$,
		\[
			\sum_{n=0}^{\infty}\left(\tfrac{3}{2}\right)^{\!n} = \infty,
		\]
		porque la razón es mayor que uno y los términos crecen en lugar de achicarse.
		\textbf{Inestable.}

		\emph{$h_4$:} la suma recorre $n\le 0$, y el cambio de variable $m=-n$ la reordena
		como una serie geométrica de razón $1/2$,
		\[
			\sum_{n=-\infty}^{0} 2^{n} = \sum_{m=0}^{\infty} 2^{-m}
			= \sum_{m=0}^{\infty}\left(\tfrac{1}{2}\right)^{\!m} = 2 < \infty.
		\]
		\textbf{Estable.} La exponencial de base $2$ crece hacia la derecha, pero como está
		encendida hacia la izquierda, lo que se recorre al sumar es su cola decreciente.

		\emph{$h_5$:} los términos no tienden a cero, así que la serie no puede converger.
		Se ve sobre un período de la sinusoide, que con $\lambda=\pi/4$ ocupa ocho muestras.
		Los módulos en esas ocho muestras son
		\[
			1,\ \tfrac{\sqrt{2}}{2},\ 0,\ \tfrac{\sqrt{2}}{2},\ 1,\ \tfrac{\sqrt{2}}{2},\ 0,\ \tfrac{\sqrt{2}}{2},
		\]
		cuya suma vale $2+2\sqrt{2}\approx 4{,}83$. Ese mismo aporte se repite en cada
		período y hay infinitos períodos, así que la suma total supera cualquier cota.
		\textbf{Inestable.} Este es el caso de estabilidad marginal: $h_5[n]$ está acotada en
		módulo por uno, pero no es absolutamente sumable, y el criterio BIBO exige lo
		segundo.

		\item La clasificación solo mira si el soporte es finito:
		$h_1$ (una muestra) y $h_2$ (cinco muestras) son \textbf{FIR}; $h_3$, $h_4$ y $h_5$
		tienen infinitas muestras no nulas y son \textbf{IIR}. Las dos primeras confirman de
		paso la regla general, ya que los dos sistemas FIR salieron estables sin necesidad
		de examinar sus coeficientes.

		\item El caso es $h_4[n]=2^n u[-n]$, el único de los
		cinco que es estable y no causal a la vez. Las dos propiedades son independientes porque
		responden preguntas distintas sobre la misma secuencia. La causalidad pregunta
		\emph{dónde} vive $h[n]$, si a la izquierda o a la derecha del origen, y no le
		importa cuánto valen sus muestras. La estabilidad pregunta \emph{cuánto suman} esas
		muestras en módulo, y no le importa dónde están ubicadas. Los ejemplos del ejercicio
		cubren las cuatro combinaciones posibles, y la Tabla~\ref{tab:sol_ej6} las reúne:
		$h_4$ es estable sin ser causal, $h_3$ y $h_5$ son causales sin ser estables, y
		$h_1$ y $h_2$ cumplen las dos.
	\end{enumerate}

	\begin{table}[H]
		\centering
		\caption{Clasificación de las cinco respuestas al impulso del ejercicio 6.}
		\label{tab:sol_ej6}
		\begin{tabular}{lccccc}
			\toprule
			$h[n]$ & Soporte & Memoria & Causal & Estable & FIR/IIR \\
			\midrule
			$3\,\delta[n]$                 & $\{0\}$           & no & sí & sí & FIR \\
			$u[n]-u[n-5]$                  & $[0,4]$           & sí & sí & sí & FIR \\
			$(3/2)^n\,u[n]$                & $[0,\infty)$      & sí & sí & no & IIR \\
			$2^n\,u[-n]$                   & $(-\infty,0]$     & sí & no & sí & IIR \\
			$\cos(\pi n/4)\,u[n]$          & $[0,\infty)$      & sí & sí & no & IIR \\
			\bottomrule
		\end{tabular}
	\end{table}

	% ============================================================================
	\tema{Ecuaciones en diferencias y realización}
	% ============================================================================

	Una \textbf{ecuación en diferencias lineal con coeficientes constantes} liga la salida
	con la entrada y con valores pasados de ambas,
	\[
		\sum_{k=0}^{N} a_k\,y[n-k] = \sum_{k=0}^{M} b_k\,x[n-k],
		\qquad a_0 \neq 0,
	\]
	y con la condición de \textbf{reposo inicial} ---todas las muestras de salida
	anteriores al comienzo de la entrada valen cero--- define un sistema LIT y causal.
	Despejando el término actual de la salida queda la \textbf{recursión}
	\[
		y[n] = \frac{1}{a_0}\left[\,\sum_{k=0}^{M} b_k\,x[n-k]
		       \;-\; \sum_{k=1}^{N} a_k\,y[n-k]\,\right],
	\]
	que se corre muestra por muestra, porque todos los valores del lado derecho ya se
	conocen cuando el cálculo llega al instante $n$. Alimentándola con $\delta[n]$ se
	obtiene $h[n]$ sin resolver ninguna ecuación.

	Una ecuación \textbf{no recursiva} ---sin términos $a_k$ con $k\ge 1$--- produce un
	sistema FIR, cuya respuesta al impulso es la lista de los coeficientes $b_k/a_0$. Una
	ecuación \textbf{recursiva} produce, en general, un sistema IIR.

	La realización usa tres bloques elementales: el \textbf{retardo unitario}
	$w[n]=v[n-1]$, el \textbf{escalador} por una constante y el \textbf{sumador}. La
	entrada atraviesa una cadena de retardos cuyos nodos se escalan por $b_k/a_0$, la
	salida atraviesa otra cadena realimentada cuyos nodos se escalan por $-a_k/a_0$, y un
	sumador reúne todas las ramas.

	% ----------------------------------------------------------------------------
	\ejercicio{7: una ecuación recursiva de primer orden}

	\begin{consigna}
	Un sistema en reposo inicial ($y[-1]=0$) responde a la ecuación en
	diferencias
	\[
		y[n] = x[n] + 0{,}8\,y[n-1].
	\]
	\begin{enumerate}
		\item Correr la recursión con $x[n]=\delta[n]$ y escribir $h[n]$
		en forma cerrada.
		\item Clasificar el sistema como FIR o IIR y determinar si es
		estable.
		\item Calcular las primeras cinco muestras de la salida ante
		$x[n]=u[n]$ y determinar hacia qué valor tiende.
		\item Graficar la realización del sistema con retardos unitarios,
		sumadores y escaladores.
		\item Determinar qué valores del coeficiente que multiplica a
		$y[n-1]$ producen un sistema inestable.
	\end{enumerate}
	\end{consigna}

	\begin{enumerate}
		\item La ecuación ya está despejada, así que se la
		corre directamente. El reposo inicial fija $y[-1]=0$, y la entrada es el impulso, que
		vale $1$ en $n=0$ y cero en el resto:
		\[
			\begin{aligned}
			y[0] &= \delta[0] + 0{,}8\cdot y[-1] = 1 + 0 = 1, \\
			y[1] &= \delta[1] + 0{,}8\cdot y[0]  = 0 + 0{,}8 = 0{,}8, \\
			y[2] &= 0 + 0{,}8\cdot 0{,}8 = 0{,}64 = (0{,}8)^2, \\
			y[3] &= 0 + 0{,}8\cdot (0{,}8)^2 = (0{,}8)^3.
			\end{aligned}
		\]
		Desde $n=1$ la entrada aporta cero y cada valor se obtiene multiplicando el anterior
		por $0{,}8$, siempre el mismo factor. El patrón se sostiene solo: si $y[n-1]=(0{,}8)^{n-1}$,
		la propia ecuación produce $y[n]=0{,}8\cdot(0{,}8)^{n-1}=(0{,}8)^n$, y como $y[0]=1$
		la igualdad vale para todo $n\ge 0$. Antes del origen la salida es nula por el reposo
		inicial, de modo que
		\[
			\boxed{\;h[n] = (0{,}8)^{n}\,u[n].\;}
		\]

		\item La exponencial no se anula para ningún
		$n \ge 0$, así que el soporte es infinito y el sistema es \textbf{IIR}. Para la
		estabilidad se suman los módulos, que forman una serie geométrica de razón
		$|0{,}8|<1$:
		\[
			\sum_{n=0}^{\infty} |0{,}8|^{n} = \frac{1}{1-0{,}8} = 5 < \infty.
		\]
		El sistema es \textbf{estable}.

		\item Ahora la entrada aporta un $1$ en cada
		instante, no solo en el origen. Corriendo la misma recursión con $y[-1]=0$:
		\[
			\begin{aligned}
			y[0] &= 1 + 0{,}8\cdot 0        = 1, \\
			y[1] &= 1 + 0{,}8\cdot 1        = 1{,}8, \\
			y[2] &= 1 + 0{,}8\cdot 1{,}8    = 2{,}44, \\
			y[3] &= 1 + 0{,}8\cdot 2{,}44   = 2{,}952, \\
			y[4] &= 1 + 0{,}8\cdot 2{,}952  = 3{,}3616.
			\end{aligned}
		\]
		La salida crece, pero cada incremento es menor que el anterior. Restando dos
		valores consecutivos de la recursión, el aporte de la entrada se cancela y queda
		$y[n]-y[n-1] = 0{,}8\,\bigl(y[n-1]-y[n-2]\bigr)$, de modo que cada incremento es
		$0{,}8$ veces el anterior ---$0{,}8$, después $0{,}64$, después $0{,}512$--- y los
		aportes se apagan geométricamente. El límite se obtiene del resultado del
		ejercicio~4, ya que la respuesta al escalón acumula la respuesta al impulso:
		\[
			\lim_{n\to\infty} y[n] = \sum_{n=0}^{\infty} (0{,}8)^{n} = \frac{1}{1-0{,}8} = 5.
		\]
		La salida tiende a $\boxed{5}$, y las cinco muestras calculadas van en camino a ese
		valor. Vale notar que el límite coincide con la suma que decidió la estabilidad en el
		inciso anterior, que es la misma relación que el ejercicio~4 estableció en general.

		\item La ecuación se transcribe directamente al diagrama de
		la Figura~\ref{fig:sol_ej7}, sin ninguna manipulación algebraica. La entrada $x[n]$
		ingresa al sumador con coeficiente $1$, de modo que no lleva escalador. La salida se
		toma a la derecha del sumador y se deriva hacia un retardo unitario, cuya salida es
		$y[n-1]$; esa muestra se escala por $0{,}8$ y vuelve al sumador. El lazo cerrado es
		lo que vuelve recursivo al sistema, y el único elemento de memoria es ese retardo.

		\begin{figure}[H]
			\centering
			\includestandalone[width=0.52\linewidth]{figures/fig_sol_ej7/fig_sol_ej7}
			\caption{Ejercicio~7. Realización de $y[n]=x[n]+0{,}8\,y[n-1]$ con un retardo
			unitario realimentado con ganancia $0{,}8$, un escalador y un sumador.}
			\label{fig:sol_ej7}
		\end{figure}

		\item Con un coeficiente genérico $a$ en lugar
		de $0{,}8$, la misma recursión produce $h[n]=a^n u[n]$ y la condición de estabilidad
		es la convergencia de la serie geométrica $\sum_{n\ge 0}|a|^n$, que converge
		únicamente cuando $|a|<1$. Por lo tanto
		\[
			\boxed{\;\text{el sistema es inestable para } |a| \ge 1.\;}
		\]
		El caso de borde merece atención. Con $|a|=1$ la respuesta al impulso no crece
		---vale $1$ para todo $n\ge0$ si $a=1$, y alterna entre $1$ y $-1$ si $a=-1$---, pero
		tampoco es absolutamente sumable, porque suma infinitos términos de módulo uno. Con
		$a=1$ la ecuación es $y[n]=y[n-1]+x[n]$, el acumulador, que ante la entrada acotada
		$x[n]=u[n]$ produce una salida que crece sin cota. Que $h[n]$ esté acotada no alcanza
		para la estabilidad BIBO.
	\end{enumerate}

	% ----------------------------------------------------------------------------
	\ejercicio{8: una ecuación recursiva de segundo orden}

	\begin{consigna}
	Un sistema en reposo inicial ($y[-1]=y[-2]=0$) responde a la ecuación en
	diferencias
	\[
		y[n] = x[n] + 0{,}9\,y[n-1] - 0{,}2\,y[n-2].
	\]
	\begin{enumerate}
		\item Correr la recursión con $x[n]=\delta[n]$ y escribir las
		primeras seis muestras de $h[n]$.
		\item Clasificar el sistema como FIR o IIR e indicar cuántas
		condiciones iniciales necesita la recursión para arrancar.
		\item Calcular las primeras cinco muestras de la salida ante
		$x[n]=u[n]$ y determinar hacia qué valor tiende, sin resolver la
		recursión en forma cerrada.
		\item Graficar la realización del sistema con retardos unitarios,
		sumadores y escaladores.
		\item Rehacer las primeras seis muestras de $h[n]$ con el
		coeficiente de $y[n-2]$ cambiado de $-0{,}2$ a $+0{,}2$, indicar
		qué le pasa a la salida y qué responde en ese caso el
		procedimiento del inciso c) sobre el valor al que tiende.
	\end{enumerate}
	\end{consigna}

	\begin{enumerate}
		\item La ecuación ya está despejada y se la corre
		como en el ejercicio anterior, con la diferencia de que cada muestra necesita las dos
		anteriores en lugar de una sola. El reposo inicial fija $y[-1]=y[-2]=0$ y la entrada
		es el impulso, que vale $1$ en $n=0$ y cero en el resto,
		\[
			\begin{aligned}
			h[0] &= 1 + 0{,}9\cdot 0 - 0{,}2\cdot 0 = 1, \\
			h[1] &= 0 + 0{,}9\cdot 1 - 0{,}2\cdot 0 = 0{,}9, \\
			h[2] &= 0{,}9\cdot 0{,}9 - 0{,}2\cdot 1 = 0{,}81 - 0{,}2 = 0{,}61, \\
			h[3] &= 0{,}9\cdot 0{,}61 - 0{,}2\cdot 0{,}9 = 0{,}549 - 0{,}18 = 0{,}369, \\
			h[4] &= 0{,}9\cdot 0{,}369 - 0{,}2\cdot 0{,}61 = 0{,}3321 - 0{,}122 = 0{,}2101, \\
			h[5] &= 0{,}9\cdot 0{,}2101 - 0{,}2\cdot 0{,}369 = 0{,}18909 - 0{,}0738 = 0{,}11529,
			\end{aligned}
		\]
		de modo que las primeras seis muestras son
		\[
			\boxed{\;h[n] = \{\, \underset{\uparrow}{1} \quad 0{,}9 \quad 0{,}61 \quad
			0{,}369 \quad 0{,}2101 \quad 0{,}11529 \quad \ldots \,\}\;}
		\]
		Desde $n=1$ la entrada aporta cero y cada valor sale de los dos anteriores. En el
		sistema de primer orden del ejercicio~7 el cociente entre una muestra y la anterior
		valía siempre $0{,}8$, y ahí se reconocía la potencia $(0{,}8)^n$. Acá ese cociente
		cambia en cada paso, $0{,}9$ primero y después $0{,}678$, $0{,}605$ y $0{,}569$, así
		que ninguna potencia sola reproduce la secuencia. Por eso el enunciado pide las
		muestras y no la forma cerrada.

		\item La recursión arranca en $n=0$,
		y para armar esa primera muestra la ecuación pide $y[-1]$ e $y[-2]$. Necesita
		entonces \textbf{dos condiciones iniciales}, una por cada muestra pasada de la salida
		que aparece en la ecuación.

		Para el soporte alcanza con mirar qué necesitaría la respuesta para apagarse. Pasado
		$n=0$ la entrada ya no aporta nada y cada muestra queda armada con las dos
		anteriores, de modo que la única forma de que la secuencia se anule de ahí en
		adelante es que dos muestras consecutivas valgan cero a la vez. Las calculadas
		decrecen sin llegar a cero, así que el soporte es infinito y el sistema es
		\textbf{IIR}.

		\item Ahora la entrada vale $1$ en todo instante
		$n\ge 0$, no solo en el origen. Corriendo la misma recursión con $y[-1]=y[-2]=0$,
		\[
			\begin{aligned}
			y[0] &= 1 + 0{,}9\cdot 0 - 0{,}2\cdot 0        = 1, \\
			y[1] &= 1 + 0{,}9\cdot 1 - 0{,}2\cdot 0        = 1{,}9, \\
			y[2] &= 1 + 0{,}9\cdot 1{,}9 - 0{,}2\cdot 1    = 2{,}51, \\
			y[3] &= 1 + 0{,}9\cdot 2{,}51 - 0{,}2\cdot 1{,}9 = 2{,}879, \\
			y[4] &= 1 + 0{,}9\cdot 2{,}879 - 0{,}2\cdot 2{,}51 = 3{,}0891.
			\end{aligned}
		\]
		La salida crece con incrementos cada vez menores, $0{,}9$ entre las dos primeras
		muestras y después $0{,}61$, $0{,}369$ y $0{,}2101$. Esos incrementos son exactamente
		las muestras de $h[n]$ del inciso a), que es la relación $g[n]-g[n-1]=h[n]$ del
		ejercicio~4 apareciendo sobre los números.

		El valor al que tiende se obtiene sin resolver la recursión. Si la salida se acerca a
		un valor $Y$ al aumentar $n$, entonces $y[n]$, $y[n-1]$ e $y[n-2]$ se acercan las tres
		a ese mismo $Y$, y como la entrada vale $1$ la ecuación pasa a ser
		\[
			Y = 1 + 0{,}9\,Y - 0{,}2\,Y.
		\]
		Agrupando los términos en $Y$ queda $Y\,(1-0{,}9+0{,}2)=1$, es decir $0{,}3\,Y=1$ y
		\[
			\boxed{\;\lim_{n\to\infty} y[n] = \frac{1}{0{,}3} = \frac{10}{3} \approx 3{,}333.\;}
		\]
		Las cinco muestras calculadas van en camino a ese valor, y la comprobación sale de
		acumular la respuesta al impulso, porque $1+0{,}9+0{,}61+0{,}369+0{,}2101+0{,}11529 =
		3{,}2044$ ya está cerca del total. La cuenta supone que el límite existe, y el
		inciso~e) muestra qué ocurre cuando esa suposición no se cumple.

		\item La ecuación se transcribe al diagrama de la
		Figura~\ref{fig:sol_ej8} sin ninguna manipulación algebraica. La entrada $x[n]$ llega
		al sumador de arriba con coeficiente $1$ y no lleva escalador. La salida se toma a la
		derecha de ese sumador y se deriva hacia una cascada de dos retardos unitarios, que
		producen $y[n-1]$ y $y[n-2]$. Cada una se escala por su coeficiente, $0{,}9$ y
		$-0{,}2$, y el sumador de abajo junta esos dos aportes antes de que el de arriba les
		sume la entrada. Comparado con el sistema de primer orden del ejercicio~7, el lazo
		tiene ahora dos elementos de memoria y dos escaladores, que es justamente lo que
		obliga a las dos condiciones iniciales del inciso b).

		\begin{figure}[H]
			\centering
			\includestandalone[width=0.62\linewidth]{figures/fig_sol_ej8/fig_sol_ej8}
			\caption{Ejercicio~8. Realización de $y[n]=x[n]+0{,}9\,y[n-1]-0{,}2\,y[n-2]$ con
			dos retardos unitarios en cascada, un escalador por cada muestra pasada de la
			salida y dos sumadores.}
			\label{fig:sol_ej8}
		\end{figure}

		\item Con $+0{,}2$ en lugar de $-0{,}2$ la
		ecuación es $y[n]=x[n]+0{,}9\,y[n-1]+0{,}2\,y[n-2]$, y la recursión se corre igual que
		antes,
		\[
			\begin{aligned}
			h[0] &= 1, \\
			h[1] &= 0{,}9\cdot 1 + 0{,}2\cdot 0 = 0{,}9, \\
			h[2] &= 0{,}9\cdot 0{,}9 + 0{,}2\cdot 1 = 0{,}81 + 0{,}2 = 1{,}01, \\
			h[3] &= 0{,}9\cdot 1{,}01 + 0{,}2\cdot 0{,}9 = 0{,}909 + 0{,}18 = 1{,}089, \\
			h[4] &= 0{,}9\cdot 1{,}089 + 0{,}2\cdot 1{,}01 = 0{,}9801 + 0{,}202 = 1{,}1821, \\
			h[5] &= 0{,}9\cdot 1{,}1821 + 0{,}2\cdot 1{,}089 = 1{,}06389 + 0{,}2178 = 1{,}28169.
			\end{aligned}
		\]
		Las dos primeras muestras coinciden con las del sistema original, porque $h[1]$
		todavía no usa el coeficiente que cambió. Desde $h[2]$ los dos aportes se suman en
		vez de restarse y la secuencia crece en lugar de decrecer. La
		Figura~\ref{fig:sol_ej8e} pone las dos respuestas al impulso una al lado de la otra,
		en la misma escala, para ver el contraste sobre las mismas seis muestras.

		\begin{figure}[H]
			\centering
			\includestandalone[width=0.78\linewidth]{figures/fig_sol_ej8e/fig_sol_ej8e}
			\caption{Ejercicio~8e. (a) La respuesta al impulso con el coeficiente $-0{,}2$,
			que decrece muestra a muestra. (b) La misma recursión con $+0{,}2$, que crece.
			Los dos paneles comparten la escala vertical.}
			\label{fig:sol_ej8e}
		\end{figure}

		Aplicado a esta ecuación, el procedimiento del inciso c) responde algo imposible. El
		planteo del valor final es ahora $Y = 1 + 0{,}9\,Y + 0{,}2\,Y$, que agrupado queda
		$Y\,(1-1{,}1)=1$ y despeja $Y=-10$. La respuesta al escalón de este sistema acumula
		muestras positivas y en $n=0$ ya vale $1$, así que nunca podría acercarse a un número
		negativo. El absurdo señala dónde falla el planteo, porque la cuenta del inciso c)
		supone que la salida tiende a algún valor y este sistema no cumple esa condición.

		Ese mismo razonamiento decide la estabilidad sin necesidad de una forma cerrada. Si
		la respuesta al impulso fuera absolutamente sumable, la respuesta al escalón tendería
		al total acumulado, que es un número positivo por ser suma de muestras positivas, y
		ese total tendría que satisfacer la ecuación del valor final, que solo admite $-10$.
		Como un mismo número no puede ser positivo y valer $-10$, la suma de los módulos
		diverge y
		\[
			\boxed{\;\text{el sistema con } +0{,}2 \text{ es inestable.}\;}
		\]
	\end{enumerate}

	% ----------------------------------------------------------------------------
	\ejercicio{9: dos ecuaciones para el mismo sistema}

	\begin{consigna}
	Sean los sistemas
	\begin{align*}
		\mathcal{T}_A&: \; y[n] = x[n]+x[n-1]+x[n-2]+x[n-3], \\
		\mathcal{T}_B&: \; y[n] = y[n-1] + x[n] - x[n-4].
	\end{align*}
	\begin{enumerate}
		\item Obtener $h[n]$ de cada uno corriendo la recursión con
		$x[n]=\delta[n]$, tomando $y[-1]=0$ en $\mathcal{T}_B$.
		\item Indicar si los dos sistemas son el mismo y justificar la
		respuesta sobre las dos respuestas al impulso.
		\item Graficar la realización de cada ecuación e indicar cuántos
		retardos y cuántas sumas usa cada una.
		\item Generalizar las dos ecuaciones a una ventana de $L$
		muestras, contar las sumas que cada forma hace por muestra de
		salida e indicar cuál conviene implementar cuando $L$ es grande.
	\end{enumerate}
	\end{consigna}

	\begin{enumerate}
		\item \emph{$\mathcal{T}_A$.} La ecuación no es recursiva, así que la respuesta al impulso
		es la lista de los coeficientes. Reemplazando la entrada por $\delta[n]$,
		\[
			h_A[n] = \delta[n]+\delta[n-1]+\delta[n-2]+\delta[n-3]
			       = \{\, \underset{\uparrow}{1} \quad 1 \quad 1 \quad 1 \,\}.
		\]

		\emph{$\mathcal{T}_B$.} Aquí hay que correr la recursión, con reposo inicial
		$y[-1]=0$ y recordando que $\delta[n]$ vale uno solo en $n=0$:
		\[
			\begin{aligned}
			y[0] &= 0 + \delta[0] - \delta[-4] = 0 + 1 - 0 = 1, \\
			y[1] &= 1 + \delta[1] - \delta[-3] = 1 + 0 - 0 = 1, \\
			y[2] &= 1 + 0 - 0 = 1, \\
			y[3] &= 1 + 0 - 0 = 1, \\
			y[4] &= 1 + \delta[4] - \delta[0] = 1 + 0 - 1 = 0, \\
			y[5] &= 0 + 0 - \delta[1] = 0.
			\end{aligned}
		\]
		Desde $n=5$ los tres términos son nulos y la salida se mantiene en cero, de modo que
		\[
			\boxed{\;h_A[n] = h_B[n] = \{\, \underset{\uparrow}{1} \quad 1 \quad 1 \quad 1 \,\}.\;}
		\]
		El paso decisivo ocurre en $n=4$, donde el término $-x[n-4]$ resta el impulso que
		había ingresado en el origen y apaga la salida para siempre. La ecuación de
		$\mathcal{T}_B$ es recursiva y su respuesta al impulso ocupa apenas cuatro muestras,
		así que el sistema es FIR. La regla general dice que una ecuación recursiva produce
		\emph{en general} un sistema IIR, y este es uno de los casos que quedan fuera de ese
		``en general''.

		\item Las dos ecuaciones definen sistemas LIT, y la
		respuesta al impulso caracteriza por completo a un sistema LIT: dos sistemas con la
		misma $h[n]$ producen la misma salida ante toda entrada, porque las dos salidas se
		calculan con la misma convolución $y[n]=x[n]*h[n]$. Como $h_A[n]=h_B[n]$, los dos
		sistemas son indistinguibles mirando sus entradas y salidas.

		Lo que ambos calculan se lee en $h[n]$: cuatro pesos iguales a uno sobre las cuatro
		últimas muestras, es decir la suma móvil de la ventana actual,
		\[
			y[n] = \sum_{k=0}^{3} x[n-k].
		\]
		$\mathcal{T}_A$ la escribe sumando las cuatro muestras cada vez, y $\mathcal{T}_B$ la
		escribe partiendo del total anterior, agregándole la muestra que ingresa y quitándole
		la que sale de la ventana. Es la misma cuenta organizada de dos maneras.

		\item Están en la Figura~\ref{fig:sol_ej9}, y cada
		una transcribe su propia ecuación.

		En $\mathcal{T}_A$ la entrada atraviesa una cadena de \textbf{tres retardos}, y los
		cuatro nodos de esa cadena ---$x[n]$, $x[n-1]$, $x[n-2]$ y $x[n-3]$--- se reúnen en
		la línea de sumadores. Reunir cuatro términos exige \textbf{tres sumas}. No hay
		escaladores, porque los cuatro coeficientes valen uno.

		En $\mathcal{T}_B$ la entrada atraviesa una cadena de \textbf{cuatro retardos} para
		producir $x[n-4]$, y la salida atraviesa \textbf{un retardo} más en el lazo
		realimentado, con lo que son cinco en total. El sumador reúne tres términos, $x[n]$
		con signo más, $x[n-4]$ con signo menos y $y[n-1]$ con signo más, lo que exige
		\textbf{dos sumas}. Tampoco lleva escaladores.

		La forma recursiva usa entonces más memoria y menos aritmética, que es el
		intercambio que el inciso siguiente cuantifica.

		\begin{figure}[H]
			\centering
			\includestandalone[width=0.86\linewidth]{figures/fig_sol_ej9/fig_sol_ej9}
			\caption{Ejercicio~9. (a) Realización no recursiva de $\mathcal{T}_A$, con tres
			retardos sobre la entrada y tres sumas. (b) Realización recursiva de
			$\mathcal{T}_B$, con cuatro retardos sobre la entrada, uno sobre la salida
			realimentada y dos sumas.}
			\label{fig:sol_ej9}
		\end{figure}

		\item Las dos ecuaciones se generalizan sin cambiar de
		forma. La no recursiva suma las $L$ últimas muestras y la recursiva actualiza el
		total anterior,
		\[
			\mathcal{T}_A: \; y[n] = \sum_{k=0}^{L-1} x[n-k],
			\qquad
			\mathcal{T}_B: \; y[n] = y[n-1] + x[n] - x[n-L].
		\]
		El conteo por muestra de salida está en la Tabla~\ref{tab:sol_ej8}. En
		$\mathcal{T}_A$ hay que reunir $L$ términos, lo que exige $L-1$ sumas, y esa cantidad
		crece con la ventana. En $\mathcal{T}_B$ hay que reunir tres términos siempre, lo que
		exige \textbf{dos sumas cualquiera sea $L$}.

		\begin{table}[H]
			\centering
			\caption{Costo de las dos formas para una ventana de $L$ muestras.}
			\label{tab:sol_ej8}
			\begin{tabular}{lcc}
				\toprule
				 & $\mathcal{T}_A$ (no recursiva) & $\mathcal{T}_B$ (recursiva) \\
				\midrule
				Retardos                & $L-1$ & $L+1$ \\
				Sumas por muestra       & $L-1$ & $2$   \\
				\bottomrule
			\end{tabular}
		\end{table}

		Con $L$ grande conviene la forma \textbf{recursiva}. La cantidad de retardos es
		prácticamente la misma en las dos ---difieren en dos posiciones de memoria, sin
		importar cuánto valga $L$---, mientras que el trabajo aritmético por muestra pasa de
		crecer proporcionalmente a $L$ a mantenerse constante. Para una ventana de mil
		muestras, $\mathcal{T}_A$ hace $999$ sumas por cada valor de la salida y
		$\mathcal{T}_B$ hace dos.

		El enunciado compara las dos formas en aritmética exacta. En una implementación real,
		donde cada valor se guarda con una cantidad finita de bits, la forma recursiva tiene
		una contrapartida: como arrastra el total de un paso al siguiente, los errores de
		redondeo se acumulan en ese total, mientras que la forma no recursiva recalcula la
		suma completa cada vez y no los propaga.
	\end{enumerate}

\end{document}
